% =====================================================================
% passages_rapport.tex
% Fragment a coller dans rapport.tex (pas de preambule ici).
% Paquets necessaires : babel (french), graphicx, booktabs, array.
% Figures attendues dans le dossier du rapport :
%   confusion_dtw_final.png, exp03_locuteur_echafaudage.png,
%   exp05_references_echafaudage.png, fig_benchmark_precision_parametres.png
% =====================================================================

\newcommand{\fichemodele}[6]{%
  \begin{center}\footnotesize
  \renewcommand{\arraystretch}{1.15}
  \begin{tabular}{@{}>{\raggedright\arraybackslash}p{0.20\linewidth}>{\raggedright\arraybackslash}p{0.76\linewidth}@{}}
    \toprule
    \textbf{Référence} & #1\\
    \midrule
    \textbf{Idée} & #2\\
    \textbf{Protocole et résultats} & #3\\
    \textbf{Coût} & #4\\
    \textbf{Limites} & #5\\
    \textbf{Utilité pour nous} & #6\\
    \bottomrule
  \end{tabular}
  \end{center}}

% =====================================================================
\section{Résultats des expériences sur la DTW et les HMM}

Sauf mention contraire, les précisions sont mesurées sur 1\,000 mots de test dits par des locuteurs absents des références, avec un intervalle de confiance (IC) à 95~\% de $\pm$\,3,1 points. Un écart inférieur à cette valeur ne se commente pas. Le test compte 10 mots, d'où un hasard à 10~\%. « Échafaudage » désigne la première chaîne fournie par l'énoncé (26 colonnes par trame, fonction \texttt{dtw} de MATLAB). « Version finale » désigne nos caractéristiques (39 colonnes) avec notre fonction \texttt{maDTW}.

\subsection{Exp01 : précision sur signal propre}

Avec MFCC, $\Delta$, $\Delta\Delta$ et CMVN et 20 références par mot, la DTW atteint 57,2~\% (IC 95~\% : $\pm$\,3,1 points), pour un F1 macro de 0,567. L'échafaudage donnait 57,6~\%. L'écart de 0,4 point est très inférieur à l'incertitude.

La matrice de confusion (figure~\ref{fig:confusion}) sépare trois groupes. « yes » (79~\%), « stop » (74~\%), « no » (70~\%) et « on » (69~\%) sont les mieux reconnus. « left » (43~\%), « off » (39~\%) et « up » (35~\%) sont les moins bien reconnus. Les plus grosses confusions sont up pris pour off (24 fois sur 100), off pris pour on (22 fois), puis off pris pour up, down pris pour no, go pris pour no et left pris pour yes (20 fois chacune).

Nous n'avons pas cherché la cause de ces confusions. Pour no et go, seule la consonne initiale diffère. Pour off et up, la consonne finale est brève et peu énergique, ce qui peut la rendre difficile à distinguer ; c'est une hypothèse que nous n'avons pas vérifiée.

Le temps moyen est de 103,4 ms par mot, extraction comprise, contre 11,7 ms pour l'échafaudage. La version finale cumule deux changements, \texttt{maDTW} (écrite avec des boucles) et 39 colonnes au lieu de 26, et nous n'avons pas séparé leurs effets.

\begin{figure}[ht]
  \centering
  \includegraphics[width=0.85\linewidth]{confusion_dtw_final.png}
  \caption{Matrice de confusion de la DTW (20 références par mot, 100 fichiers de test par mot, précision 57,2~\%). Les lignes donnent le mot dit, les colonnes le mot reconnu.}
  \label{fig:confusion}
\end{figure}

\subsection{Exp02 : bruit}

Le test est en conditions non appariées : les références sont propres et le mot à reconnaître est bruité. Nous ajoutons un bruit blanc ou un bruit de robinet à 20, 10, 5 et 0~dB de SNR, calculé sur tout le clip, silences compris. Ces mesures datent de l'échafaudage (57,6~\% sur signal propre).

\begin{table}[ht]
  \centering
  \begin{tabular}{lcc}
    \toprule
    SNR & Bruit blanc & Bruit de robinet\\
    \midrule
    Signal propre & 57,6~\% & 57,6~\%\\
    20~dB & 56,7~\% & 54,9~\%\\
    10~dB & 49,1~\% & 40,2~\%\\
    5~dB & 17,0~\% & 15,1~\%\\
    0~dB & 13,3~\% & 14,0~\%\\
    \bottomrule
  \end{tabular}
  \caption{Précision de la DTW avec références propres et test bruité (échafaudage, IC 95~\% d'environ $\pm$\,3 points).}
\end{table}

La précision tient jusqu'à 20~dB, baisse à 10~dB et tombe à 5 et 0~dB près du hasard (10~\%). Le résultat va à l'encontre de notre attente. Le bruit blanc couvre toutes les fréquences, et nous pensions qu'il serait le plus destructeur pour les MFCC. À 10~dB, c'est le bruit de robinet qui dégrade le plus (40,2~\% contre 49,1~\%). Aux autres niveaux, l'écart entre les deux bruits reste dans l'incertitude.

Nous proposons deux explications, sans les avoir testées. Le bruit de robinet est irrégulier et son énergie occupe des bandes proches de celles de la parole, alors que le bruit blanc, constant, abîme toutes les trames de la même façon et la DTW compare alors des mots abîmés de façon comparable. Ou bien la détection de parole est trompée : un bruit irrégulier ressemble davantage à de la parole et peut faire couper le mot au mauvais endroit.

\subsection{Exp03 : locuteurs connus et inconnus}

La DTW atteint 57,2~\% ($\pm$\,3,1 points, 1\,000 mots) sur des locuteurs jamais vus et 83,5~\% ($\pm$\,3,3 points, 491 mots) sur des fichiers des locuteurs des références, soit 26 points d'écart (échafaudage : 57,6~\% et 83,9~\%). Les deux intervalles sont loin de se toucher (figure~\ref{fig:locuteur}). La méthode mémorise des exemples de voix et dépend donc fortement du locuteur. L'écart peut aussi refléter des conditions d'enregistrement proches entre la référence et le test d'un même locuteur (micro, pièce) ; nous ne séparons pas les deux effets.

\begin{figure}[ht]
  \centering
  \includegraphics[width=0.7\linewidth]{exp03_locuteur_echafaudage.png}
  \caption{Précision de la DTW selon que le locuteur est connu ou inconnu (échafaudage). Les barres d'erreur sont les IC à 95~\% : $\pm$\,3,1 points sur 1\,000 mots pour les locuteurs inconnus, $\pm$\,3,3 points sur 491 mots pour les locuteurs connus.}
  \label{fig:locuteur}
\end{figure}

\subsection{Exp04 : choix des caractéristiques}

On garde le même système et on change seulement la description du son.

\begin{table}[ht]
  \centering
  \begin{tabular}{lcc}
    \toprule
    Variante & Colonnes & Précision\\
    \midrule
    FBANK & 26 & 58,9~\%\\
    MFCC & 13 & 58,0~\%\\
    MFCC + $\Delta$ + $\Delta\Delta$ & 39 & 58,6~\%\\
    MFCC + $\Delta$ + $\Delta\Delta$ + CMVN & 39 & 57,2~\%\\
    idem, fenêtre de Hann & 39 & 56,6~\%\\
    \bottomrule
  \end{tabular}
  \caption{Précision de la DTW selon les caractéristiques (1\,000 mots, IC 95~\% : $\pm$\,3,1 points).}
\end{table}

Les cinq précisions tiennent dans 2,3 points (56,6 à 58,9~\%), pour une incertitude de $\pm$\,3,1 points. Aucune différence n'est statistiquement significative : nous ne pouvons conclure ni que FBANK vaut mieux que MFCC, ni que la CMVN fait perdre 1,4 point, ni que Hann vaut moins que Hamming. L'ordre annoncé par le guide (FBANK $<$ MFCC $<$ MFCC + deltas $\le$ CMVN) n'apparaît pas, et notre prévision que la CMVN ferait le plus de différence est fausse. La variante 4 redonne 57,2~\%, comme l'exp01 : le script de comparaison et la chaîne utilisée ailleurs concordent.

Trois explications sont possibles, aucune n'est testée. Le signal est propre, et la CMVN corrige surtout les différences de micro et de canal. La DTW capte déjà la dynamique temporelle, donc les deltas ajoutent peu. Avec 20 références par mot, la limite vient peut-être de la variabilité entre locuteurs plutôt que de la représentation.

Une conséquence pratique : les MFCC seuls (13 colonnes) stockeraient environ 146\,000 valeurs au lieu de 438\,633, pour une précision indiscernable de celle des autres variantes. Tout est mesuré sur signal propre. Nous n'avons pas refait la comparaison sous bruit, là où la CMVN devrait se justifier.

\subsection{Exp05 : nombre de références par mot}

La précision monte avec le nombre de références, avec des rendements décroissants : 32,4~\% avec 1 référence, 47,1~\% avec 5, 53,9~\% avec 10 et 57,6~\% avec 20 (échafaudage). Le temps de comparaison est proportionnel au nombre de références : de 0,4 à 7,8 ms par mot, soit environ 0,04 ms par comparaison. Passer de 10 à 20 références double le temps pour 3,7 points de gain, du même ordre que l'incertitude ($\pm$\,3 points). Les premières références sont les plus rentables.

\begin{figure}[ht]
  \centering
  \includegraphics[width=0.7\linewidth]{exp05_references_echafaudage.png}
  \caption{Précision et temps de reconnaissance selon le nombre de références par mot (échafaudage).}
\end{figure}

\subsection{Exp06 : coût en temps et en mémoire}

Cette expérience ne mesure pas la précision. Elle mesure ce que coûte le système, sur un Lenovo ThinkPad T14 Gen 2i (Intel Core i7-1185G7, 16~Gio, Ubuntu 26.04.1).

\begin{table}[ht]
  \centering
  \begin{tabular}{lcc}
    \toprule
     & Échafaudage & Version finale\\
    \midrule
    Précision sur signal propre & 57,6~\% & 57,2~\%\\
    F1 macro & 0,569 & 0,567\\
    Test dépendant & 83,9~\% & 83,5~\%\\
    Colonnes par trame & 26 & 39\\
    Temps moyen par mot (exp01) & 11,7 ms & 103,4 ms\\
    Extraction (exp06) & 1,7 à 1,9 ms & 0,8 à 1,0 ms\\
    Reconnaissance, un clip (exp06) & 6,8 à 7,4 ms & 58 à 61 ms\\
    Valeurs en mémoire & 193\,336 & 438\,633\\
    \bottomrule
  \end{tabular}
  \caption{Échafaudage et version finale.}
\end{table}

Pour un clip d'environ 1~s, l'échafaudage met 1,8 ms à extraire les caractéristiques et 7 ms à les comparer aux 200 références, soit un facteur temps réel de 0,009. La comparaison représente environ 80~\% du temps et croît avec le nombre de références. La réponse arrive environ 9 ms après la fin du mot. La version finale est une dizaine de fois plus lente (58 à 61 ms de comparaison), soit un facteur temps réel de l'ordre de 0,06, toujours compatible avec le temps réel.

Une réserve sur les temps : la même fonction \texttt{maDTW} est chronométrée à 58--61 ms (exp06), 70 ms (exp07), 96,7 ms (exp05) et 103,4 ms (exp01, extraction comprise). Nous n'avons pas identifié l'origine de ces écarts. Nous citons chaque temps avec l'expérience qui le donne.

La DTW n'a aucun paramètre appris. Elle garde en mémoire ses valeurs de référence : 193\,336 pour l'échafaudage (environ 1,5 Mo en double précision) et 438\,633 pour la version finale (environ 3,5 Mo). C'est ce dernier nombre qui sert dans la figure du benchmark.

\subsection{Exp07 : corrélation et bande de Sakoe-Chiba}

Cette expérience porte sur 200 mots de test seulement (IC 95~\% : $\pm$\,6,9 points), avec \texttt{maDTW} et des caractéristiques précalculées. Les temps sont donc des temps de reconnaissance seule. La corrélation normalisée compare les mots après rééchantillonnage linéaire à une même durée, sans alignement non linéaire. La demi-largeur $R$ de la bande est en trames de 10 ms : $R = 5$ vaut 50 ms.

\begin{table}[ht]
  \centering
  \begin{tabular}{lcc}
    \toprule
    Méthode & Précision & Temps par mot\\
    \midrule
    Corrélation normalisée & 49,5~\% & non mesuré\\
    DTW, $R = \infty$ & 53,5~\% & 70 ms\\
    DTW, $R = 20$ & 53,5~\% & 55 ms\\
    DTW, $R = 10$ & 53,0~\% & 48 ms\\
    DTW, $R = 5$ & 51,5~\% & 42 ms\\
    \bottomrule
  \end{tabular}
  \caption{Corrélation et DTW avec bande de Sakoe-Chiba (200 mots de test).}
\end{table}

La DTW (53,5~\%) devance la corrélation (49,5~\%), mais l'écart de 4 points reste dans l'incertitude de $\pm$\,6,9 points. Une bande de demi-largeur 20 ne change pas la précision et réduit le temps de 21~\%. Avec $R = 5$, le temps baisse de 40~\% pour 2 points perdus, ce qui est dans l'incertitude. Le gain de temps reste modeste parce que notre implémentation calcule les distances locales sur toute la matrice : la bande ne limite que la récurrence. Enfin, le 53,5~\% de cet échantillon est inférieur au 57,2~\% de l'exp01 sur 1\,000 mots ; l'écart de 3,7 points est compatible avec le tirage des 200 mots.

\subsection{Exp08 : HMM discret et HMM-GMM}

Le HMM discret remplace chaque trame par un numéro parmi 64 (quantification par k-moyennes) et apprend une suite de numéros par mot. Le HMM-GMM garde les vecteurs de caractéristiques et modélise chaque état par un mélange de gaussiennes. Les deux apprennent sur 200 exemples par mot, tirés au hasard dans l'ensemble d'apprentissage. Le nombre d'états $N$ (4, 6 ou 8) et de gaussiennes $K$ (1, 2 ou 4) se choisit sur 300 mots de validation, jamais sur le test.

\begin{table}[ht]
  \centering
  \begin{tabular}{lccc}
    \toprule
    & $K = 1$ & $K = 2$ & $K = 4$\\
    \midrule
    $N = 4$ & 72,0~\% & 77,7~\% & 81,0~\%\\
    $N = 6$ & 73,7~\% & 80,0~\% & 85,0~\%\\
    $N = 8$ & 77,3~\% & 82,3~\% & 84,3~\%\\
    \bottomrule
  \end{tabular}
  \caption{Précision du HMM-GMM sur la validation (300 mots, IC 95~\% d'environ $\pm$\,4 points).}
\end{table}

Plus de gaussiennes aident toujours : $K = 4$ gagne de 7 à 11 points sur $K = 1$. Avec $K = 4$, $N = 6$ et $N = 8$ sont indiscernables sur 300 mots ; le script retient $N = 6$. Le meilleur $K$ est peut-être au-delà de 4, car la grille s'arrête là.

\begin{table}[ht]
  \centering
  \begin{tabular}{lcccc}
    \toprule
    Méthode & Apprentissage & Précision & IC 95~\% & Temps par mot\\
    \midrule
    DTW + MFCC & 20 références/mot & 57,2~\% & $\pm$\,3,1 & 96,7 ms (exp05)\\
    HMM discret & 200 exemples/mot & 74,7~\% & $\pm$\,2,7 & 1,0 ms\\
    HMM-GMM ($N = 6$, $K = 4$) & 200 exemples/mot & 86,4~\% & $\pm$\,2,1 & 2,8 ms\\
    \bottomrule
  \end{tabular}
  \caption{Test sur 1\,000 mots, locuteurs inconnus. Les temps des HMM excluent l'extraction des caractéristiques (environ 1 ms de plus).}
  \label{tab:hmm}
\end{table}

Le classement DTW $<$ HMM discret $<$ HMM-GMM est net : les intervalles ne se touchent pas. Le HMM-GMM garde les vecteurs réels et apprend, pour chaque état, une forme moyenne à quatre gaussiennes. Le HMM discret perd du détail en ramenant chaque trame à un seul numéro. Pour le HMM discret, la matrice de confusion montre off pris pour up (30 fois), up pris pour off (21 fois), right pris pour left (12 fois), puis left pris pour yes, go pris pour no, down pris pour go et down pris pour no (10 fois chacun). « yes » (92~\%), « no » (85~\%) et « stop » (83~\%) sont les mieux reconnus, « up » (54~\%) et « off » (60~\%) les moins bien. Nous n'avons pas la matrice du HMM-GMM, dont les F1 les plus bas sont ceux de « go » (0,75), « up » et « down » (0,81) et « off » (0,83) : les mots difficiles sont les mêmes.

Le HMM reconnaît un mot en 1 à 3 ms contre près de 100 ms pour la DTW, mais il demande un apprentissage long (11 modèles appris pour le réglage et le test). La DTW est l'inverse : aucun apprentissage, une comparaison avec toutes les références.

Deux réserves. Le script a signalé des non-convergences : \texttt{fitgmdist} pour un cas ($N = 8$, $K = 4$) et \texttt{kmeans} et \texttt{hmmtrain} pour plusieurs mots du HMM discret, qui est donc peut-être sous-entraîné. Et surtout, les HMM apprennent sur 200 exemples par mot, la DTW sur 20. Nous avions prévu de relancer les HMM avec 20 exemples par mot pour séparer l'effet du modèle de celui des données. Ce test n'a pas été mené à terme : notre script contenait une erreur que nous n'avons pas corrigée à temps. L'écart de 29 points entre la DTW et le HMM-GMM ne peut donc pas être attribué au seul modèle. Nous n'avons pas non plus refait le test locuteur connu/inconnu avec les HMM.

% =====================================================================
\section{Vérification des briques de base}

\subsection{Alignement DTW (\texttt{maDTW})}

Sur les séquences jouets du guide, notre fonction donne DTW(A, B) = 2 et DTW(A, C) = 8, comme attendu. Le chemin A-B retrouvé est (1,1) (1,2) (2,3) (3,4) (3,5) (4,6), identique à celui du guide. L'écart avec la fonction \texttt{dtw} de MATLAB est nul (0,00).

\subsection{Détection d'activité vocale (\texttt{verifierVAD})}

Nous avons comparé les limites de parole trouvées par notre fonction \texttt{vad.m} à celles de la fonction \texttt{detectSpeech} de MATLAB, sur 20 fichiers de Speech Commands (deux par mot).

L'écart absolu moyen est de 92 ms au début de la parole et de 96 ms à la fin. Ces moyennes sont tirées par quelques fichiers : les médianes sont de 45 ms au début et de 55 ms à la fin. Quinze fichiers sur 20 ont un début à moins de 60 ms de celui de \texttt{detectSpeech}, et douze sur 20 une fin à moins de 60 ms. Trois fichiers dépassent 200 ms d'écart à chaque extrémité, jusqu'à 470 ms au début (\texttt{on/b83c1acf\_nohash\_0}) et 515 ms à la fin (\texttt{right/cfbedff9\_nohash\_0}).

La figure du mot « yes » montre l'énergie qui monte pendant la prononciation et qui reste faible dans le silence et le bruit de fond. Le seuil d'énergie sépare la zone de parole du reste. \texttt{detectSpeech} n'est pas une vérité de terrain : l'écart mesure un désaccord entre deux méthodes, pas une erreur de la nôtre. Pour les fichiers très écartés, nous n'avons pas examiné le signal. Un bruit ou une énergie irrégulière dans le fichier est une cause possible.

\subsection{Banc de filtres Mel et MFCC (\texttt{verifierMFCC})}

Nous avons comparé notre banc de filtres (\texttt{bancFiltresMel.m}) à celui de \texttt{designAuditoryFilterBank}. Les deux sont formés de filtres triangulaires qui se chevauchent et s'élargissent avec la fréquence. L'écart maximal entre les deux matrices est de 0,248 : les bancs ne sont pas identiques. Une normalisation différente des filtres (par la hauteur ou par l'aire) pourrait l'expliquer, ce que la différence d'amplitude observée sur c1 irait dans le sens de ; nous ne l'avons pas vérifié.

Nous avons ensuite comparé les coefficients de \texttt{mesMFCC.m} à ceux de la fonction \texttt{mfcc} de MATLAB. Les évolutions temporelles de c1 sont très proches, avec une amplitude différente. Les corrélations sont de 1,000 pour c0, 0,999 pour c1, 0,998 pour c2, 0,997 pour c3 et 0,990 pour c4, toutes au-dessus du seuil de 0,95 fixé par le guide. Notre implémentation reproduit donc l'évolution des coefficients de référence, à un facteur d'échelle près.

% =====================================================================
\section{Benchmark des méthodes modernes}

\subsection{Méthode}

Nous n'avons exécuté aucune des méthodes de cette section. Les chiffres viennent des articles originaux, relevés dans le document 06 du projet, avec leur base, leur protocole et leurs conditions (modèle de langue, données de pré-entraînement). Chaque tableau porte sur une seule tâche, et nos mesures sont dans un tableau séparé de celui des résultats publiés. Nous n'avons pas lancé l'exemple MathWorks « Speech Command Recognition Using Deep Learning » sur nos fichiers de test, qui aurait donné un point de comparaison avec le même protocole que la DTW et les HMM.

\subsection{Fiches de lecture}

\subsubsection*{DNN-HMM}
\fichemodele
{G.~Hinton et al., « Deep Neural Networks for Acoustic Modeling in Speech Recognition: The Shared Views of Four Research Groups », \emph{IEEE Signal Processing Magazine}, 2012.}
{On garde le HMM de notre volet classique (états, transitions, Viterbi). Les GMM sont remplacés par un réseau profond qui reçoit plusieurs trames (par exemple 11 trames de FBANK) et prédit la probabilité de chaque état. La règle de Bayes convertit cette probabilité en vraisemblance.}
{TIMIT : taux d'erreur de phonèmes de 27,3~\% pour un HMM classique à états contextuels, 20,7~\% pour le DNN-HMM. Switchboard (309~h) : WER de 27,4~\% pour le meilleur GMM-HMM, 18,5~\% pour le DNN-HMM, soit 33~\% de réduction relative.}
{De l'ordre du million de paramètres. 309~h de parole transcrite pour Switchboard.}
{Il faut d'abord un HMM-GMM pour produire les alignements d'apprentissage. Les hypothèses de Markov et d'alignement par Viterbi sont conservées.}
{Seul le modèle d'émission change par rapport à notre HMM-GMM (86,4~\% sur nos 10 mots). C'est le lien direct entre le volet classique et le volet avancé.}

\subsubsection*{TDNN (Kaldi)}
\fichemodele
{V.~Peddinti, D.~Povey, S.~Khudanpur, « A time delay neural network architecture for efficient modeling of long temporal contexts », Interspeech 2015. Modèle d'origine : Waibel et al., 1989.}
{Chaque couche combine des trames à des décalages temporels choisis (par exemple $t-2$, $t$, $t+2$), de plus en plus espacés dans les couches profondes. C'est une convolution 1D en temps avec sous-échantillonnage.}
{Plusieurs tâches de grand vocabulaire, de 3 à 1\,800~h de données : amélioration relative moyenne de 5,5~\% par rapport au DNN de base, et de 7,3~\% sur Switchboard, à temps d'entraînement comparable.}
{De l'ordre du million de paramètres, efficace en calcul.}
{Utilisé avec un HMM. La boîte à outils Kaldi n'est pas accessible depuis MATLAB : nous ne pouvons que citer les résultats publiés.}
{Bon compromis entre précision et coût de calcul. Ligne intermédiaire entre le HMM-GMM et les modèles de bout en bout.}

\subsubsection*{Conformer}
\fichemodele
{A.~Gulati et al., « Conformer: Convolution-augmented Transformer for Speech Recognition », Interspeech 2020, arXiv:2005.08100.}
{Chaque bloc combine de l'auto-attention (dépendances lointaines) et des convolutions (motifs locaux). Modèle de bout en bout : FBANK en entrée, texte en sortie, sans HMM.}
{LibriSpeech, WER (test-clean / test-other) sans modèle de langue : 2,7 / 6,3 pour S, 2,3 / 5,0 pour M, 2,1 / 4,3 pour L. Avec modèle de langue : 2,1 / 5,0, 2,0 / 4,3 et 1,9 / 3,9.}
{10,3 M (S), 30,7 M (M) et 118,8 M (L) de paramètres. 960~h de parole transcrite.}
{Demande beaucoup de données transcrites. Peu interprétable.}
{La version S (10,3 M) atteint 2,7~\% sur test-clean, comme Whisper large-v2 (1\,550 M), mais reste en retrait sur test-other (6,3~\% contre 5,2~\%). Les deux ne sont pas entraînés dans les mêmes conditions (voir Whisper).}

\subsubsection*{wav2vec 2.0}
\fichemodele
{A.~Baevski, H.~Zhou, A.~Mohamed, M.~Auli, « wav2vec 2.0: A Framework for Self-Supervised Learning of Speech Representations », NeurIPS 2020, arXiv:2006.11477.}
{Entrée : le signal brut. Un CNN produit une représentation toutes les 20 ms ; une partie est masquée et un Transformer doit retrouver la bonne représentation quantifiée parmi plusieurs candidats. Pré-entraînement sur de l'audio non transcrit, puis affinage sur peu de données transcrites.}
{LibriSpeech, 960~h transcrites : WER 1,8 / 3,3 (LARGE, avec modèle de langue Transformer). Avec 10 minutes transcrites et 53\,000~h non transcrites : WER 4,8 / 8,2. SUPERB (modèle gelé, Speech Commands v1, 12 classes) : 96,23~\% (BASE) et 96,66~\% (LARGE).}
{95 M (BASE) et 317 M (LARGE) de paramètres. 53\,000~h d'audio non transcrit pour le pré-entraînement.}
{Calcul lourd, pas embarquable tel quel. Le pré-entraînement est en anglais.}
{Argument pour les langues peu dotées (section Perspectives). La condition à rappeler : les 10 minutes transcrites reposent sur 53\,000~h d'audio non transcrit.}

\subsubsection*{HuBERT}
\fichemodele
{W.-N.~Hsu et al., « HuBERT: Self-Supervised Speech Representation Learning by Masked Prediction of Hidden Units », \emph{IEEE/ACM Transactions on Audio, Speech, and Language Processing}, 2021, arXiv:2106.07447.}
{Même principe de masquage que wav2vec 2.0, mais le réseau prédit l'étiquette de groupe (cluster) des zones masquées. À la première itération, les groupes viennent de k-moyennes sur des MFCC. Les itérations suivantes utilisent les représentations du modèle.}
{X-LARGE affiné sur 960~h : WER 1,8 / 2,9 (pré-entraînement 60\,000~h, modèle de langue Transformer). SUPERB, Speech Commands v1, 12 classes : 96,30~\% (BASE) et 95,29~\% (LARGE).}
{95 M (BASE), 317 M (LARGE) et 964 M (X-LARGE) de paramètres.}
{Mêmes limites que wav2vec 2.0 : taille, calcul, pré-entraînement en anglais.}
{La quantification en groupes par k-moyennes sur MFCC est la même opération que celle de notre HMM discret.}

\subsubsection*{Whisper}
\fichemodele
{A.~Radford et al., « Robust Speech Recognition via Large-Scale Weak Supervision », ICML 2023, arXiv:2212.04356.}
{Transformer encodeur-décodeur entraîné sur 680\,000~h d'audio transcrit récupéré sur le web, dans de nombreuses langues. Les transcriptions sont imparfaites (supervision faible). Entrée : spectrogramme log-Mel. Sortie : texte. Utilisé tel quel, sans affinage sur la base de test.}
{LibriSpeech test-clean, décodage glouton, sans affinage : WER de 7,6~\% (tiny), 5,0~\% (base), 3,4~\% (small), 2,9~\% (medium) et 2,7~\% (large-v2, 5,2~\% sur test-other).}
{De 39 M (tiny) à 1\,550 M (large-v2) de paramètres. 680\,000~h d'entraînement.}
{Conçu pour des segments de 30 s, pas pour des commandes d'une seconde. Sa sortie est du texte libre (« Yes. », « yeah ») qu'il faut normaliser pour la comparer à une classe. Très gros pour l'embarqué. Nécessite Python.}
{Référence pour la généralisation hors domaine (bruit, accents). Faire transcrire nos fichiers de test et compter les bonnes réponses serait un test possible ; nous ne l'avons pas fait.}

\subsubsection*{Reconnaissance embarquée (TinyML)}
\fichemodele
{P.~Warden, arXiv:1804.03209, 2018 ; Y.~Zhang et al., « Hello Edge: Keyword Spotting on Microcontrollers », arXiv:1711.07128, 2017 ; B.~Kim et al., « Broadcasted Residual Learning for Efficient Keyword Spotting », Interspeech 2021, arXiv:2106.04140 ; A.~Berg et al., « Keyword Transformer », Interspeech 2021, arXiv:2104.00769.}
{Détection de mots-clés sur microcontrôleur : petits CNN à convolutions séparables en profondeur, quantification des poids et des activations sur 8 bits (mémoire divisée par 4 par rapport au flottant 32 bits), MFCC ou log-Mel calculés sur la puce. Outil de référence : TensorFlow Lite for Microcontrollers, avec l'exemple « micro\_speech » (yes et no).}
{Speech Commands, 12 classes : voir le tableau~\ref{tab:sc}. De 94,4~\% (DS-CNN S, 38,6 Ko) à 98,7~\% (BC-ResNet-8, 321 k paramètres).}
{De 9,2 k paramètres (BC-ResNet-1) à 5,4 M (KWT-3). DS-CNN S tient dans 38,6 Ko.}
{Les chiffres mélangent les versions v1 et v2 de la base. Les verrous sont les fausses alarmes et l'ajout de mots sans réentraîner.}
{C'est la seule tâche identique à la nôtre. Elle donne la comparaison de mémoire avec la DTW (438\,633 valeurs) et les HMM.}

\subsection{Résultats publiés par tâche}

\begin{table}[ht]
  \centering\footnotesize
  \begin{tabular}{llccl}
    \toprule
    Modèle & Version & Précision & Taille & Source\\
    \midrule
    CNN de référence & v2 & 88,2~\% & non précisée & Warden, 2018\\
    DS-CNN S & v1 & 94,4~\% & 38,6 Ko (8 bits) & Zhang et al., 2017\\
    DS-CNN L & v1 & 95,4~\% & 497,6 Ko (8 bits) & Zhang et al., 2017\\
    BC-ResNet-1 & v2 & 96,9~\% & 9,2 k param. & Kim et al., 2021\\
    BC-ResNet-8 & v2 & 98,7~\% & 321 k param. & Kim et al., 2021\\
    KWT-1 & v2 & 97,72~\% & 607 k param. & Berg et al., 2021\\
    KWT-3 (distillation) & v2 & 98,56~\% & 5\,361 k param. & Berg et al., 2021\\
    wav2vec 2.0 BASE (gelé) & v1 & 96,23~\% & 95 M param. & Yang et al., 2021 (SUPERB)\\
    wav2vec 2.0 LARGE (gelé) & v1 & 96,66~\% & 317 M param. & Yang et al., 2021\\
    HuBERT BASE (gelé) & v1 & 96,30~\% & 95 M param. & Yang et al., 2021\\
    HuBERT LARGE (gelé) & v1 & 95,29~\% & 317 M param. & Yang et al., 2021\\
    \bottomrule
  \end{tabular}
  \caption{Résultats publiés sur Speech Commands, protocole à 12 classes. Dans SUPERB, le modèle est gelé et seul un petit classifieur est entraîné ; les FBANK seules y obtiennent 8,63~\%, ce qui mesure ce classifieur simple et non les FBANK en général.}
  \label{tab:sc}
\end{table}

\begin{table}[ht]
  \centering\footnotesize
  \begin{tabular}{lcccl}
    \toprule
    Modèle & Paramètres & test-clean & test-other & Conditions\\
    \midrule
    Whisper tiny & 39 M & 7,6 & n.r. & sans entraînement sur LibriSpeech\\
    Whisper base & 74 M & 5,0 & n.r. & idem\\
    Whisper small & 244 M & 3,4 & n.r. & idem\\
    Whisper medium & 769 M & 2,9 & n.r. & idem\\
    Whisper large-v2 & 1\,550 M & 2,7 & 5,2 & idem\\
    Conformer S & 10,3 M & 2,7 & 6,3 & sans modèle de langue\\
    Conformer S & 10,3 M & 2,1 & 5,0 & avec modèle de langue\\
    Conformer M & 30,7 M & 2,3 & 5,0 & sans modèle de langue\\
    Conformer M & 30,7 M & 2,0 & 4,3 & avec modèle de langue\\
    Conformer L & 118,8 M & 2,1 & 4,3 & sans modèle de langue\\
    Conformer L & 118,8 M & 1,9 & 3,9 & avec modèle de langue LSTM\\
    wav2vec 2.0 LARGE & 317 M & 1,8 & 3,3 & 53\,000~h non transcrites + modèle de langue\\
    HuBERT X-LARGE & 964 M & 1,8 & 2,9 & 60\,000~h non transcrites + modèle de langue\\
    \bottomrule
  \end{tabular}
  \caption{WER (\%) publiés sur LibriSpeech (960~h). n.r. : non relevé dans notre source. Les conditions diffèrent d'une ligne à l'autre, on ne compare pas directement un modèle sans entraînement sur LibriSpeech et un modèle affiné avec modèle de langue.}
\end{table}

\begin{table}[ht]
  \centering\footnotesize
  \begin{tabular}{lccl}
    \toprule
    Tâche (Hinton et al., 2012) & GMM-HMM & DNN-HMM & Métrique\\
    \midrule
    TIMIT (phonèmes) & 27,3~\% & 20,7~\% & taux d'erreur de phonèmes\\
    Switchboard (309~h) & 27,4~\% & 18,5~\% & WER\\
    \bottomrule
  \end{tabular}
  \caption{Avant et après les réseaux profonds. Pour le TDNN, Peddinti et al. (2015) rapportent un gain relatif de 7,3~\% sur Switchboard par rapport au DNN de base.}
\end{table}

\subsection{Nos mesures et les résultats publiés}

Le tableau~\ref{tab:hmm} donne nos mesures. La figure~\ref{fig:benchmark} les place à côté des résultats publiés sur Speech Commands. Pour la DTW, le nombre en abscisse est celui des valeurs stockées (438\,633, mesuré). Pour les HMM, c'est une estimation de l'ordre de grandeur : environ 13\,000 pour le HMM-GMM (6 états, 4 gaussiennes à covariance diagonale, 26 colonnes, 10 mots) et environ 5\,000 pour le HMM discret.

\begin{figure}[ht]
  \centering
  \includegraphics[width=\linewidth]{fig_benchmark_precision_parametres.png}
  \caption{Précision en fonction du nombre de paramètres. Les points publiés portent sur 12 classes (les dix mots, plus « silence » et « inconnu »), avec les versions v1 ou v2 de la base ; nos points portent sur 10 mots et 1\,000 fichiers de test. Les deux séries ne sont pas mesurées sur le même test.}
  \label{fig:benchmark}
\end{figure}

La comparaison directe est impossible : les 12 classes publiées incluent « silence » et « inconnu », nos 10 mots non, et la version de la base et le partage ne sont pas identiques. Trois observations tiennent malgré cela. D'abord, les modèles publiés dépassent 96~\% à partir de 9\,200 paramètres (BC-ResNet-1), alors que notre HMM-GMM, du même ordre de taille, atteint 86,4~\%. Ensuite, passer de 9\,200 à 321\,000 paramètres apporte 1,8 point (96,9~\% à 98,7~\%), et passer à 5,4 M avec un Transformer n'apporte rien de plus (98,56~\%). Enfin, notre DTW stocke 438\,633 valeurs, soit environ 3,5 Mo en double précision et 440 Ko si on les quantifie sur 8 bits, de l'ordre du DS-CNN L (497,6 Ko), pour 57,2~\%.

Sur la donnée, wav2vec 2.0 atteint un WER de 4,8 / 8,2 avec 10 minutes de transcription, mais après un pré-entraînement sur 53\,000~h d'audio non transcrit. Sur la métrique, le WER compte les erreurs de mots sur des phrases, la précision compte les mots isolés bien reconnus. Si chaque clip contient un mot et que le système rend un mot, le WER vaut 1 moins la précision. Il peut dépasser cette valeur si le système ajoute des mots, comme Whisper et son texte libre.

\subsection{Analyse comparative}

\begin{table}[ht]
  \centering\footnotesize
  \renewcommand{\arraystretch}{1.2}
  \begin{tabular}{>{\raggedright\arraybackslash}p{0.17\linewidth}>{\raggedright\arraybackslash}p{0.22\linewidth}>{\raggedright\arraybackslash}p{0.18\linewidth}>{\raggedright\arraybackslash}p{0.17\linewidth}>{\raggedright\arraybackslash}p{0.17\linewidth}}
    \toprule
    Méthode & Hypothèses et interprétabilité & Données transcrites & Paramètres & Mise en œuvre, embarqué\\
    \midrule
    DTW & quasi-stationnarité, alignement monotone ; interprétable & quelques exemples par mot (20 ici) & références stockées (438\,633 valeurs ici) & facile sous MATLAB ; possible pour un petit vocabulaire\\
    HMM-GMM & Markov, trames indépendantes, gaussiennes ; interprétable & centaines par mot (200 ici) & dizaines de milliers (environ 13\,000 ici) & moyenne ; possible\\
    DNN-HMM, TDNN & Markov, émissions sans hypothèse gaussienne ; peu interprétable & beaucoup & millions & difficile ; possible en version réduite\\
    Conformer & aucune hypothèse explicite ; peu interprétable & beaucoup (960~h) & 10 à 120 M & difficile ; difficile\\
    wav2vec 2.0, HuBERT, Whisper & aucune hypothèse explicite ; peu interprétable & très peu après pré-entraînement, ou aucune (Whisper) & 95 M à 1,5 G & Python ; non\\
    \bottomrule
  \end{tabular}
  \caption{Comparaison des méthodes. Les valeurs « ici » sont celles de nos expériences.}
\end{table}

% =====================================================================
\section{Perspectives d'applications}

\subsection{Commandes vocales dans les langues locales du Sénégal}

\textbf{Domaine.} Agriculture (commander une pompe d'irrigation, consulter la météo), domotique et services publics, pour des utilisateurs qui parlent wolof, pulaar ou sérère et qui ne lisent pas forcément le français. \textbf{Contraintes.} Bruit extérieur, téléphones bas de gamme, connexion intermittente (donc traitement local), accents variés. \textbf{Méthode.} Ces langues ont très peu de données transcrites. La DTW ne demande que quelques exemples par mot : dans nos mesures, elle donne 32,4~\% avec une référence par mot et 57,6~\% avec 20 sur des locuteurs inconnus, mais 83,5~\% sur les locuteurs des références. Elle convient donc à un utilisateur qui enregistre sa propre voix, pas à un système ouvert à tous. Pour un système ouvert, l'apprentissage auto-supervisé obtient un WER de 4,8 / 8,2 avec 10 minutes transcrites, à condition de disposer d'audio non transcrit dans la langue ou d'un modèle pré-entraîné multilingue. \textbf{Maturité et verrous.} Les méthodes sont mûres pour l'anglais. Le verrou est la collecte de données dans ces langues, avec le consentement des locuteurs, et la tenue au bruit réel.

\subsection{Commande mains libres en environnement industriel bruité}

\textbf{Domaine.} Un technicien aux mains occupées commande une machine ou consulte une procédure par la voix. \textbf{Contraintes.} SNR souvent compris entre 0 et 10~dB, temps réel, sécurité. \textbf{Méthode.} Nos courbes situent la limite de la DTW avec MFCC : 56,7~\% (bruit blanc) et 54,9~\% (robinet) à 20~dB, mais 49,1~\% et 40,2~\% à 10~dB, puis 17,0~\% et 15,1~\% à 5~dB et 13,3~\% et 14,0~\% à 0~dB, près des 10~\% du hasard. Cette méthode seule est donc inutilisable dans la plage de bruit visée. Les pistes sont l'apprentissage avec des données bruitées (non testé ici), des modèles dont les auteurs annoncent la robustesse au bruit, comme Whisper, et un micro proche de la bouche ou un réseau de micros pour relever le SNR en amont. \textbf{Maturité et verrous.} Les bruits de machine sont irréguliers, et notre exp02 montre que le bruit irrégulier (robinet) fait plus de dégâts à 10~dB que le bruit blanc.

\subsection{Assistance aux personnes à mobilité réduite}

\textbf{Domaine.} Accessibilité : commander un fauteuil roulant, un lit médicalisé ou un appel d'urgence par quelques mots. \textbf{Contraintes.} Fiabilité très élevée (une commande mal reconnue peut être dangereuse), voix parfois atypiques (maladie, fatigue), fonctionnement hors ligne, exigences réglementaires des dispositifs médicaux, confidentialité des données vocales. \textbf{Méthode.} Un petit vocabulaire et un seul utilisateur favorisent la reconnaissance dépendante du locuteur. Notre exp03 mesure 83,5~\% pour des locuteurs connus contre 57,2~\% pour des locuteurs inconnus. Mais 83,5~\% veut dire une commande sur six mal reconnue : c'est insuffisant pour cet usage. La métrique principale devient le taux de fausses alarmes sur le silence et les mots inconnus, que notre test ne mesure pas. \textbf{Maturité et verrous.} Atteindre la fiabilité requise, adapter le système à une voix atypique et passer les exigences réglementaires.

% =====================================================================
\section{Conclusion}

Sur dix mots de Speech Commands, la DTW avec 20 références par mot atteint 57,2~\% sur des locuteurs inconnus et 83,5~\% sur les locuteurs des références. Ces 26 points d'écart montrent qu'une méthode qui mémorise des voix est limitée par la variabilité des locuteurs. Sa précision monte avec le nombre de références, par gains décroissants, et le temps de calcul suit. Elle tient jusqu'à 20~dB de SNR et tombe sous 10~dB. Sur signal propre, les cinq choix de caractéristiques testés donnent des résultats indiscernables. Sa mémoire (438\,633 valeurs) dépasse de loin celle de petits modèles publiés.

Les HMM atteignent 74,7~\% (discret) et 86,4~\% (GMM), avec dix fois plus d'exemples par mot que la DTW. Nous n'avons pas pu séparer la part du modèle de celle des données. Les modèles publiés sur Speech Commands dépassent 96~\% sur 12 classes, avec 9\,200 paramètres pour le plus petit, mais sur un test qui n'est pas le nôtre.

Quatre limites restent. Le test à 20 exemples par mot pour les HMM n'a pas abouti. Le test locuteur connu/inconnu n'a pas été refait avec les HMM. La CMVN n'a pas été comparée sous bruit. L'exemple MathWorks n'a pas été exécuté, ce qui nous prive d'une comparaison sur les mêmes fichiers.
